% TOP_PROBLEM: {"id": 4700001, "title": "Low degree rigid systems", "category_id": 11, "category": {"id": 11, "name": "dynamical_systems", "display_name": "Dynamical Systems", "description": "Problems about long-term behavior of deterministic systems, Hamiltonian dynamics, and periodic orbits.", "slug": "dynamical-systems", "order_index": 11, "created_at": "2026-07-31T15:26:25.671Z"}, "status": "open", "problem_number": "AMR-046-0001", "set_id": 13}
% FIRST_POSTED: 2026-10-02
% ENTRY_KIND: statement_only
% UnsolvedMath ID 4700001; source catalogue code AMR-046-0001
% !TeX program = lualatex
% Definition and sources only; this file is not an LLM solution attempt.
\documentclass[11pt,a4paper]{article}
\usepackage[margin=25mm]{geometry}
\usepackage{fontspec}
\setmainfont{lmroman10-regular.otf}[BoldFont=lmroman10-bold.otf,
 ItalicFont=lmroman10-italic.otf,BoldItalicFont=lmroman10-bolditalic.otf]
\usepackage{amsmath,amssymb,mathtools}
\usepackage{unicode-math}
\setmathfont{latinmodern-math.otf}
\AtBeginDocument{\let\setminus\smallsetminus}
\usepackage{xurl,hyperref}
\hypersetup{colorlinks=true,urlcolor=blue}
\setcounter{secnumdepth}{0}
\setlength{\parindent}{0pt}
\setlength{\parskip}{5pt}
\setlength{\emergencystretch}{3em}
\begin{document}
\section*{Low degree rigid systems}\label{4700001}
{\small UnsolvedMath ID 4700001 · AMR-046-0001 · Dynamical Systems · AMR Open Problem Lists}
\subsection{Definitions and mathematical statement}
For real parameters $a,b,c,d,e$, consider the polynomial vector field
\[
 \dot x=-y+xF(x,y),\qquad \dot y=x+yF(x,y),\qquad
 F=a+bx+cy+dx^2+exy.
\]
A limit cycle means the image of a nonconstant periodic solution which is
isolated among periodic orbits. Distinct time parametrizations of the same
closed curve count as one cycle; multiplicity is not included in the count.
For $r>0$, the substitution $(x,y)=(r\cos\theta,r\sin\theta)$ gives
$\dot\theta=1$ and
\[
 \frac{dr}{d\theta}=ar+(b\cos\theta+c\sin\theta)r^2
 +(d\cos^2\theta+e\cos\theta\sin\theta)r^3.
\]
Thus a cycle corresponds to an isolated positive $2\pi$-periodic solution of
this scalar equation, defined without blow-up throughout the period.
The question is whether every member of the family has at most two limit
cycles, and whether this bound is attained. Equivalently, is the supremum of
the number of isolated positive periodic solutions, over all five parameters,
equal to two?

The origin is the only equilibrium, and every periodic orbit surrounds it.
Consequently the count includes all sizes of cycles, including those far from
the origin, rather than only cycles appearing in a local bifurcation. If the
origin is a center, its continuous family of surrounding periodic orbits is
not counted as a family of limit cycles. The source records examples with two
cycles; the outstanding content is a uniform global upper bound.
\subsection{Short English statement}
Can any cubic rigid system in this five-parameter family have more than two isolated periodic orbits?
\subsection{Sources}
\begin{itemize}
\item[S1] ULAM.AI and the UnsolvedMath Contributors, supplied problems.json, record 4700001 (AMR-046-0001), AMR Open Problem Lists. Catalogue identity and source transcription; CC BY 4.0.
\url{https://huggingface.co/datasets/ulamai/UnsolvedMath}.
\item[S2] Gasull - Some open problems in low dimensional dynamical systems (2020), Problem environment 1: Low degree rigid systems. Original source indexed by AMR-046-0001.
\url{https://arxiv.org/abs/2012.02524}.
\end{itemize}
\end{document}
